% Standalone source-record documentation; no research attempts.
\documentclass[11pt]{article}
\usepackage[margin=25mm]{geometry}
\usepackage{fontspec,unicode-math}
\setmainfont{Latin Modern Roman}
\setmathfont{Latin Modern Math}
\usepackage{amsmath,amssymb,mathtools,xurl,hyperref}
\hypersetup{colorlinks=true,urlcolor=blue,linkcolor=blue}
\setcounter{secnumdepth}{0}
\setlength{\parindent}{0pt}
\setlength{\parskip}{5pt}
\setlength{\emergencystretch}{3em}
\newcommand{\sref}[2]{\hyperlink{source:#1:#2}{[S#2]}}
\newcommand{\cataloguescope}[1]{\paragraph{Source record.}#1}
\newcommand{\floor}[1]{\left\lfloor#1\right\rfloor}
\begin{document}
% UnsolvedMath ID 30004671; source catalogue code OWR-7155442-005
\setcounter{section}{30004670}
\section{Geometric Analogues of Finite Effective Complexity}\label{30004671}
{\small\sffamily UnsolvedMath ID 30004671 · Computer Science}
\subsection{Definitions and mathematical statement}

\paragraph{Shared notation.}Standard number systems are $\mathbb N,\mathbb Z,\mathbb Q,\mathbb R,\mathbb C$. Specialized notation and index ranges follow the cited source and the definitions below.


A projective plane has points and lines with incidence relations; a geometric inequality compares numerical invariants of a configuration. Effective complexity needs a specified quantitative description or representation measure.
\paragraph{Record scope and qualifications.}
The catalogue title is retained as metadata. Named certificates, contracts, bounds or frameworks in that title are not asserted as results of the original source extract.
The original statement, projective invariant, and sphere-area inequality are missing from the extraction.

Source-record limitation: The finite effective-complexity statement and sphere-area inequality referred to by the question are absent.
\paragraph{Statement recorded in the supplied source.}
Can one formulate and prove the geometric analogue of the finite effective-complexity statement for points and lines on the projective plane, together with the corresponding sphere-area inequality? \sref{30004671}{2}
\paragraph{Source status.}
The supplied catalogue labels this source record open; that metadata was checked against the supplied snapshot on 2026-10-11. The label is preserved for provenance and is not an independent verification of present open status or mathematical correctness. The catalogue literature assessment is dated 2026-08-22.
\subsection{Short English statement}
Develop geometric counterparts of the cited effective-complexity statement.
\subsection{Sources}
\begin{itemize}
\item[\hypertarget{source:30004671:1}{S1}] Ulam AI and the UnsolvedMath Contributors, UnsolvedMath: A Curated Collection of Open Mathematics Problems, record 30004671 (OWR-7155442-005). Catalogue curation and metadata: CC BY 4.0. The source extract is retained; definitions, English summaries and qualifications are editorial additions.
\url{https://huggingface.co/datasets/ulamai/UnsolvedMath}.
\item[\hypertarget{source:30004671:2}{S2}] Computability Theory (2022), Oberwolfach Reports, DOI 10.4171/owr/2021/21, pp. 1149--1190
\url{https://doi.org/10.4171/owr/2021/21}.
\item[\hypertarget{source:30004671:3}{S3}] Oberwolfach Report 21/2021
\url{https://ems.press/content/serial-article-files/46899}.
\item Source-specific Oberwolfach Reports licences, including CC BY-SA 4.0 where applicable, are retained; see THIRD\_PARTY\_NOTICES.md.
\end{itemize}
\paragraph{Preparation and review.}
The original source and UnsolvedMath contributors retain attribution for the source extract. Codex (GPT-6; exact serving revision not exposed) assisted with definitions, English summaries, scope notes and formatting. No new human mathematical authorship or independently verified solution is asserted. Permanent source ID 30004671 and code OWR-7155442-005 checked against the complete supplied registry and current repository. Distinct upstream records, including repeated questions and broad or damaged extracts, are retained for complete coverage of the source CS category. A repeated record is not a new independent conjecture. Complete mathematical-equivalence and current-literature review has not been performed. Preparation and final review follow contribution prompts C and D, with source gaps disclosed. Structural and rendering checks do not establish a new theorem or open status.
\label{end:30004671}
\end{document}
